\documentclass[12pt]{article}
\usepackage[utf8]{inputenc}
\usepackage[T1]{fontenc}
\usepackage{hyperref}
\title{Shabo Language Resources for K-12 Curriculum Development}
\author{Compiled by Claude Sonnet 5.5 (claude-sonnet-5-5) for Andr\'as Kornai}
\date{}
\begin{document}
\maketitle

\section{Introduction}
Shabo (also Chabu, Mikeyir, Mekeyer; ISO 639-3 sbf; Glottolog shab1252) is an endangered language isolate of south-western Ethiopia, spoken near the Sheka forest. Estimates run from about 400 to 600 speakers. Few children now acquire it as a first language, with speakers shifting to Majang and Shekkacho, and the vocabulary shows heavy borrowing from Majang and Amharic. No school materials were found in search results.

\section{Primary Educational Resources}
No Shabo primers, textbooks or school programmes were found.

\section{Linguistic and Reference Materials}
\begin{itemize}
\item ELDP collection 0074, Documentation and grammatical description of Shabo (Kibebe Tsehay, Addis Ababa University): \url{https://eldp.access.preservica.com/dk0074/}
\item Culture in Crisis project page for the same Shabo documentation project (began 2010): \url{https://cultureincrisis.org/projects/documentation-and-grammatical-description-of-the-shabo-language-a-very-endangered-isolate-language-of-southwestern-ethiopia}
\item Shabo language, Wikipedia (cites the Kibebe 2015 grammar and Teferra Anbessa's 1991 sketch): \url{https://en.wikipedia.org/wiki/Shabo_language}
\item WALS entry for Shabo: \url{https://wals.info/languoid/lect/wals_code_shb}
\item Academic paper on Shabo as an isolate (Academia.edu): \url{https://www.academia.edu/39266165}
\end{itemize}

\section{Religious and Community Resources}
No Shabo scripture or community organisation was found in search results.

\section{Digital and Audio Resources}
\begin{itemize}
\item Endangered Languages Project entry for Shabo: \url{https://endangeredlanguages.com/lang/1379?hl=en}
\item Rosetta Project / Panglossia entry for Shabo (sbf): \url{https://rosettapanglossia.longnow.org/index.php?title=Sbf}
\end{itemize}
The ELDP collection is marked online and may contain audio and video recordings; its contents were not examined here.

\section{Curriculum Development Context}
Shabo is a very small community language under shift pressure, with a grammar (2015) and an ELDP documentation archive but no confirmed orthography in use or school materials. A curriculum would most realistically be a heritage-language revitalisation component built on the ELDP recordings, working with Addis Ababa University and the community, alongside Amharic and the regionally dominant languages.

\end{document}